\documentclass[12pt]{article}
\usepackage[T1]{fontenc}
\usepackage[utf8]{inputenc}
\usepackage{lmodern}
\usepackage{textcomp}
\usepackage{enumitem}
\usepackage{geometry}
\geometry{margin=1in}
\begin{document}
\begin{center}
Answer Key - Philosophy - Graduate Level Assessment 21 \
\small (Answers are ordered exactly as in the question paper.)
\end{center}

\section*{Multiple Choice}
\begin{enumerate}[label=\arabic*.]
\item \textbf{Answer:} A
\\ \textit{Options:} A) Kannon; B) Shakyamuni; C) Manjusri; D) Guan-yin; E) Tara
\\ \textit{Note:} Kannon, also known as Avalokiteshvara, is revered in Mahayana Buddhism for her compassionate role in guiding the souls of the deceased, particularly children, towards salvation and enlightenment.
\item \textbf{Answer:} A
\\ \textit{Options:} A) Bast; B) Nehebkau; C) Wag and Thoth; D) Tsagaan Sar; E) Songkran
\\ \textit{Note:} Bast is the correct answer because it represents the significant cultural and religious aspects of ancient Babylon, particularly during their New Year celebrations, which were integral to their identity and philosophical understanding of time and renewal.
\item \textbf{Answer:} A
\\ \textit{Options:} A) the notion of "freedom" is overly simplistic.; B) the concept of "pleasure" is unclear.; C) the notion of "legislating for oneself" is absurd.; D) the idea of "moral truth" is a contradiction.; E) we can't really distinguish between good and evil.
\\ \textit{Note:} Anscombe argues that Butler's conception of freedom fails to account for the complexities of moral agency and the various influences that shape human behavior, leading to an overly simplistic understanding of what it means to be free.
\item \textbf{Answer:} D
\\ \textit{Options:} A) the fundamental laws of physics.; B) the purpose of life.; C) the existence of the divine.; D) the causes of good and evil.; E) both a and b.
\\ \textit{Note:} Augustine asserts that understanding the causes of good and evil is essential for achieving true happiness, as it enables individuals to align their desires and actions with the ultimate good, thus fostering a deeper sense of fulfillment and purpose.
\item \textbf{Answer:} A
\\ \textit{Options:} A) the same level of attention and resources as those with more native assets.; B) educational resources based on their social background.; C) resources based on their willingness to learn.; D) educational resources based on their economic background.; E) less attention and fewer resources than those with more native assets.
\\ \textit{Note:} Rawls advocates for the principle of fairness and equality of opportunity, arguing that society should ensure that individuals with fewer native assets receive equal attention and resources to help them achieve their potential and mitigate disadvantages.
\end{enumerate}

\section*{True / False}
\begin{enumerate}[label=\arabic*.]
\setcounter{enumi}{5}
\item \textbf{Answer:} True
\\ \textit{Options:} A) True; B) False
\\ \textit{Note:} Aquinas argues that the ultimate end of human life is the beatific vision, which is the direct contemplation of God, leading to perfect fulfillment and happiness.
\item \textbf{Answer:} False
\\ \textit{Options:} A) True; B) False
\\ \textit{Note:} Cicero emphasizes the importance of reason and virtue in achieving moral integrity, suggesting that blindly following societal norms can lead to moral failings rather than true ethical behavior.
\item \textbf{Answer:} True
\\ \textit{Options:} A) True; B) False
\\ \textit{Note:} Wolf argues that interpreting Kant's ethics as a finite set of constraints undermines its core tenet of moral autonomy and the idea that ethical principles must be universally applicable and grounded in rationality.
\item \textbf{Answer:} False
\\ \textit{Options:} A) True; B) False
\\ \textit{Note:} The Battle of Uhud did not directly lead to the conversion of Mecca to Islam; instead, it was the subsequent Treaty of Hudaybiyyah and the strategic political and social developments that followed that ultimately facilitated Mecca's conversion.
\item \textbf{Answer:} False
\\ \textit{Options:} A) True; B) False
\\ \textit{Note:} The statement is false because while Japan implemented a policy of isolation known as Sakoku from 1635 to 1853, it did not completely exclude all forms of Christianity and European influence during that entire period, as there were instances of limited contact and the presence of hidden Christians.
\end{enumerate}

\section*{Long Form Response}
\begin{enumerate}[label=\arabic*.]
\setcounter{enumi}{10}
\item \textbf{Answer:} The statement "If either George enrolls or Harry enrolls, then Ira does not enroll" can be logically structured as (G ∨ H) → \textasciitilde{}I, where G represents George enrolling, H represents Harry enrolling, and I represents Ira enrolling. This formulation indicates that the enrollment of either George or Harry is sufficient to conclude that Ira will not enroll. The implication of this formal representation is significant as it highlights a conditional relationship where the disjunction of George and Harry's enrollment leads to the negation of Ira's enrollment. This structure allows for the exploration of scenarios in which either George or Harry's decision directly influences Ira's choice, thus illustrating the interdependencies of their actions within a logical framework. Overall, this representation provides clarity in understanding the implications of enrollment decisions among the individuals involved.
\\ \textit{Note:} The answer correctly identifies the logical structure of the statement as a conditional involving a disjunction, demonstrating that the enrollment of either George or Harry directly leads to the conclusion that Ira will not enroll, thereby emphasizing the nature of conditional relationships in propositional logic.
\item \textbf{Answer:} Ashoka's promotion of non-violence as a form of dharma represents a significant philosophical shift that emphasizes the intrinsic value of life and ethical conduct over the traditional notions of power and conquest prevalent in his time. By adopting and propagating the principles of ahimsa, or non-harm, Ashoka redefined moral responsibility, suggesting that true power lies in compassion and restraint rather than domination. This approach not only transformed the ethical landscape of his own reign, encouraging a more humane governance model, but it also laid foundational ideas that influenced Buddhist thought and later ethical frameworks, promoting peace and coexistence. The enduring legacy of Ashoka's non-violent dharma continues to resonate in contemporary discussions on ethics, social justice, and conflict resolution, showcasing the profound impact of his philosophical stance on both historical and modern ethical thought.
\\ \textit{Note:} correct because it articulates how Ashoka's embrace of non-violence and ahimsa redefined ethical responsibility by valuing compassion and moral integrity over traditional metrics of power, thereby influencing ethical thought both in his era and in subsequent philosophical discourse.
\end{enumerate}

\vfill
\noindent\textit{End of Answer Key}
\end{document}